\documentclass{article}
\usepackage[utf8]{inputenc}
\usepackage{amsmath}
\usepackage{geometry}
\geometry{margin=2cm}
\begin{document}

\section*{Análise da Questão 120 — ENEM 2024 — Ciências da Natureza}

\subsection*{Interpretação do Enunciado, Estratégia e Análise Preliminar}
A questão apresenta uma célula a combustível microbiana (CCM) que gera eletricidade pela oxidação do acetato (CH$_3$COO$^-$) e redução de oxigênio, dividida em duas câmaras: ânodo e cátodo. A reação no ânodo mostra a oxidação do acetato com liberação de CO$_2$, íons H$^+$ e elétrons. No cátodo ocorre a redução do oxigênio pela captura de elétrons e íons H$^+$. A questão pede a reação global que representa o funcionamento completo da CCM.

\paragraph{Termos importantes}
\begin{itemize}
  \item \textbf{Reação global:} Combinação das semi-reações do ânodo e do cátodo eliminando intermediários como elétrons.
  \item \textbf{Semi-reações:} Reações químicas individuais no ânodo e cátodo.
  \item \textbf{Balanceamento:} Ajuste dos coeficientes para conservar massa e carga.
\end{itemize}

\paragraph{O que a questão pede}:
Encontrar a reação química global balanceada que integra as semi-reações do ânodo e cátodo.

\paragraph{Estratégia de resolução:}
Somar as semi-reações indicadas, multiplicando-as para igualar o número de elétrons liberados e consumidos, então somar e cancelar elétrons e íons H$^+$.

\subsection*{Conteúdo e Mini Revisão}
\begin{itemize}
\item \textbf{Células a combustível microbianas (CCM):} Utilizam bactérias eletrogênicas para oxidar matéria orgânica no ânodo, liberando elétrons e H$^+$. Estes elétrons viajam até o cátodo, onde o oxigênio é reduzido formando água.

\item \textbf{Balanceamento:} Fundamental para garantir conservação de carga e massa; elétrons liberados no ânodo devem ser consumidos no cátodo.

\item \textbf{Reações apresentadas:}
\[
\begin{aligned}
\text{Ânodo: } & \mathrm{CH_3COO^- + 2 H_2O \rightarrow 2 CO_2 + 7 H^+ + 8 e^-} \\
\text{Cátodo: } & \mathrm{4 H^+ + O_2 + 4 e^- \rightarrow 2 H_2O}
\end{aligned}
\]

\end{itemize}

\subsection*{Resolução e \'Pulo do Gato'}
\begin{enumerate}
\item Multiplicar reação do cátodo por 2 para igualar elétrons:
\[
\mathrm{8 H^+ + 2 O_2 + 8 e^- \rightarrow 4 H_2O}
\]
\item Somar as semi-reações:
\[
\mathrm{CH_3COO^- + 2 H_2O + 8 H^+ + 2 O_2 + 8 e^- \rightarrow 2 CO_2 + 7 H^+ + 8 e^- + 4 H_2O}
\]
\item Cancelar elétrons e adequar íons:
\[
\mathrm{CH_3COO^- + H^+ + 2 O_2 \rightarrow 2 CO_2 + 2 H_2O}
\]

\paragraph{Pulo do gato:} Ajustar coeficientes para igualar elétrons e somar para cancelar intermediários simplifica a reação global de forma rápida e correta.
\end{enumerate}

\subsection*{Armadilhas Comuns e Análise das Alternativas}
\begin{itemize}
\item \textbf{A:} Reação incompleta, não inclui a redução do oxigênio nem o balanceamento correto.
\item \textbf{B:} Apresenta a reação no sentido inverso, confundindo síntese com oxidação.
\item \textbf{D:} Reação incorreta e mal balanceada, não representa o funcionamento real da CCM.
\item \textbf{E:} Reação invertida e mal balanceada, oposta ao funcionamento esperado.
\end{itemize}

\subsection*{Código da Questão}
\texttt{CN\_FISIO\_QUIM\_001}

% Referência bibliográfica omitida conforme política de evitar identificação de pessoas reais ou trabalhos específicos.

\end{document}